\cgeentry{cge_gams}{Write the model as a GAMS program, and solve it in GAMS}

\cgesec{Description}

\cmd{cge_gams} writes the calibrated model and one scenario as a GAMS program
in a folder. Anyone with GAMS can then solve the model with PATH, CONOPT,
KNITRO, IPOPT or another solver, check the results of \pkg with a second
implementation, or extend the equations. GAMS receives what a user supplies,
not the parameters that Stata calibrated: the SAM after the layout rewrites
of \manref{cge_map}, the roles of its accounts and the settings given to
\manref{cge_calibrate}. The GAMS program calibrates every parameter itself,
solves the base and then the scenario. Its equations are an implementation
of the model written in GAMS, not translated from Mata.

\cmd{cge_gams} can also start GAMS from interactive Stata, compare the GAMS
results with the Stata solution, export every period of a dynamic run, and
import a GAMS solution as a Stata simulation.

\cgesec{Quick start}

Write the model and the scenario of simulation \cmd{notariff} to folder
\cmd{gams_notariff}, to solve with \cmd{gams cge_run.gms} there
\begin{stcode}
cge_gams using gams_notariff, sim(notariff)
\end{stcode}
Write it again, solve it in GAMS with PATH and compare the results with Stata's
\begin{stcode}
cge_gams using gams_notariff, sim(notariff) replace run compare gams(C:/GAMS/54/gams.exe)
\end{stcode}
Solve the files already written with CONOPT, and compare
\begin{stcode}
cge_gams using gams_notariff, sim(notariff) noexport run compare solver(conopt) ///
    gams(C:/GAMS/54/gams.exe)
\end{stcode}
Compare results that GAMS wrote earlier, for example after a batch run
\begin{stcode}
cge_gams using gams_notariff, sim(notariff) noexport compare
\end{stcode}
Export the current shocks, and after solving in GAMS store the solution as
simulation \cmd{big_path}
\begin{stcode}
cge_gams using big, replace
cge_gams using big, noexport import(big_path)
\end{stcode}
Export every period of dynamic run \cmd{base}
\begin{stcode}
cge_gams using gams_base, dynamic(base)
\end{stcode}

\cgesec{Syntax}

\begin{stsyntax}
\cmdb{cge_gams} \cmd{using} \argu{folder} [\texttt{,} \argu{options}]
\end{stsyntax}

\begin{opttable}
\optgroup{Scenario}
\optarg{}{sim}{name} & export the shocks and closure of a stored simulation\\
\optarg{dyn}{amic}{name} & export every period of a dynamic run\\
\opt{}{replace} & replace the files of an earlier export\\
\optgroup{GAMS}
\opt{}{run} & solve the exported model in GAMS\\
\optarg{}{gams}{path} & the GAMS executable; default \cmd{gams}\\
\optarg{}{solver}{name} & the solver; default \cmd{path}\\
\optarg{}{type}{name} & GAMS model type: \cmd{cns}, \cmd{nlp} or \cmd{mcp}\\
\optgroup{Results}
\opt{}{compare} & compare the GAMS results with the Stata solution\\
\optarg{imp}{ort}{name} & store the GAMS solution as simulation \argu{name}\\
\opt{}{noexport} & use the files already in \argu{folder}\\
\end{opttable}

\cmd{sim()} and \cmd{dynamic()} may not be combined, nor \cmd{import()} and
\cmd{dynamic()}. \opt{}{compare} needs \cmd{sim()} or \cmd{dynamic()}, and
\opt{}{noexport} needs \opt{}{run}, \opt{}{compare} or \cmd{import()}.

\cgesec{Options}

\begin{optdesc}
\item[\optarg{}{sim}{name}] exports the shocks and closure of a stored
simulation. Without \cmd{sim()} or \cmd{dynamic()}, the shocks and closure in
memory are exported, the scenario that the next \manref{cge_solve} would
solve, under the name \cmd{current}. A period of a dynamic run is exported
with \cmd{dynamic()}.

\item[\optarg{dyn}{amic}{name}] exports every period of a run of
\manref{cge_dynamic}. The changes of each period from the period before, to
parameters and to fixed variables, are written to
\mbox{\cmd{cge_}\argu{sim}\cmd{_par.inc}} and \mbox{\cmd{cge_}\argu{sim}\cmd{_var.inc}},
where \argu{sim} is the period's simulation, such as \cmd{base2025}, and the
driver solves the periods in order, each from the solution of the one before.
The periods must have the same closure, and runs with constraints are not
exported.

\item[\opt{}{replace}] allows the files of an earlier export in \argu{folder}
to be replaced. The folder, and any parent folder, is created if it does not
exist.

\item[\opt{}{run}] solves the exported model in GAMS. Stata deletes any
earlier result files, writes a script in the folder, \cmd{cge_run.cmd} on
Windows and \cmd{cge_run.sh} elsewhere, and runs it with \cmd{shell}; the
script can be run again by hand. It runs
\cmd{gams cge_run.gms lo=2 lf=cge_run.log} with the solver and model type
chosen. Stata for Windows cannot start other programs in batch mode: there,
\opt{}{run} is an error, and \cmd{cge_run.gms} is run with GAMS separately and
the results compared with \opt{}{noexport} \opt{}{compare}.

\item[\optarg{}{gams}{path}] gives the GAMS executable, for example
\cmd{gams(C:/GAMS/54/gams.exe)}. By default, \cmd{gams} must be on the system
path.

\item[\optarg{}{solver}{name}] chooses the solver; the default is \cmd{path}.
When GAMS runs, it writes an option file with tolerances scaled to the size
of the SAM for PATH, CONOPT, IPOPT and KNITRO; other solvers use their
defaults.

\item[\optarg{}{type}{name}] chooses the GAMS model type: \cmd{cns}, a
constrained nonlinear system (PATH, CONOPT, KNITRO); \cmd{nlp}, the same
equations with a constant objective, for NLP solvers; or \cmd{mcp}, a mixed
complementarity problem (PATH), which a scenario with constraints needs. The
default is \cmd{mcp} for a static scenario with constraints, \cmd{nlp} for
the solvers \cmd{ipopt}, \cmd{ipopth}, \cmd{minos} and \cmd{snopt}, and
\cmd{cns} otherwise; a scenario with constraints cannot be exported as
another type. Both the solver and the type can also be set when GAMS is run,
as in \texttt{gams cge\_run.gms -{}-cgesolver=ipopt -{}-cgetype=nlp}.

\item[\opt{}{compare}] compares the GAMS results with the Stata base and with
simulation \cmd{sim()}, or with every period of the run of \cmd{dynamic()}.
It reports the GAMS model status and, for the base and for the scenario or
each period, the number of values compared and the largest relative
difference, with the variable where it occurs.

\item[\optarg{imp}{ort}{name}] stores the GAMS solution of a static scenario
as simulation \argu{name}: the free variables take the values of GAMS, the
fixed variables those of Stata, and the model's equations are evaluated in
Stata at the result, whose largest scaled residual is reported. A model too
large for Stata's solver can thus be solved by PATH and its results used in
Stata with \manref{cge_report}, \manref{cge_results} or \manref{cge_microsim}.
With \opt{}{noexport}, the results already in \argu{folder} are imported.

\item[\opt{}{noexport}] uses the files already in \argu{folder}, which must
hold \cmd{cge_run.gms}: with \opt{}{run} it solves them again, for example
after they have been edited, and with \opt{}{compare} or \cmd{import()} alone
it uses results that GAMS wrote earlier.
\end{optdesc}

\cgesec{Remarks and examples}

The GAMS program is a second implementation of the model of Part~II. GAMS
calibrates the parameters from the SAM and the settings and solves the base
and the scenario itself, so a GAMS solution that matches Stata's confirms the
calibration as well as the solution. The program can also be given to a
modeller who works in GAMS, or extended with equations that the package does
not have. \Chapref{chap:gams} describes the GAMS implementation.

\cgesubsec{Example 1: exporting a scenario}

The synthetic SAM is calibrated as in \manref{cge_solve}, and every import
tariff is removed, with the sales tax rates replacing the revenue:
\stlog{gams_sim}
\cmd{cge_gams} writes the model and the scenario to folder
\cmd{gams_notariff}, which it creates:
\stlog{gams_export}
Two parameters change from the base: the tariffs on agriculture and
manufactures (the base tariff on services is zero). The changes are written as
ratios to the base values, so that GAMS applies them to its own calibrated
values, and the closure follows as codes for the GAMS program:
\stlog{gams_par}
The driver, \cmd{cge_run.gms}, includes the model, then the scenario, and
states how to run it:
\stlog{gams_driver}
An earlier export is not overwritten without \opt{}{replace}:
\stlog{gams_again}

\cgesubsec{Example 2: solving in GAMS}

In interactive Stata, \opt{}{run} starts GAMS and \opt{}{compare} reads its
results, for example
\begin{stcode}
. cge_gams using gams_notariff, sim(notariff) replace run compare gams(C:/GAMS/54/gams.exe)
\end{stcode}
Stata for Windows cannot start GAMS in batch mode, so a do-file run in batch
mode stops:
\stlog{gams_batch}
In batch mode, \cmd{gams cge_run.gms} is run in the folder outside Stata, and
\begin{stcode}
. cge_gams using gams_notariff, sim(notariff) noexport compare
\end{stcode}
then compares the results. A GAMS licence covering the solver is needed to
solve a model of realistic size; GAMS runs solvers outside the licence only
for small models. GAMS listing files (\cmd{.lst}) name the holder of the
licence.

\cgesubsec{Example 3: a scenario with constraints}

A constraint set by \manref{cge_constraint} makes the scenario a mixed
complementarity problem. Here imports of manufactures may rise by at most 2
per cent when tariffs are removed, with the premium on the quota going to the
government:
\stlog{gams_quota}
The constraint becomes a slack variable \cmd{SLK1}, zero or positive, paired
with the bound \cmd{BND1}, and the adjusting variable, the import premium
\cmd{RM(com_man)}, is freed and linked to the slack by \cmd{LNK1}. The model
type is \cmd{mcp} by default, for PATH; CONOPT, KNITRO and IPOPT do not solve
complementarity problems. From version 0.9.8, the premium starts at its fixed
level, zero, so that GAMS finds the regime itself; before, it was written at
the solution. GAMS writes the room within the bound and its marginal, the
premium, which \opt{}{compare} checks against the solution listing of
\manref{cge_listing}.

\cgesubsec{Example 4: a dynamic run}

With \cmd{dynamic()}, every period of a run is exported. A run of two periods
after the base, with capital specific to each activity:
\stlog{gams_dyn}
Each period has its own pair of scenario files, with its changes from the
period before: the growth of government flows and of the base of the health
link in the parameter file, and the stocks of capital and labour, which are
fixed variables, in the variable file:
\stlog{gams_dynfiles}

\cgesubsec{The files}

\begin{cgetable}{>{\ttfamily}l p{0.66\linewidth}}
\toprule
\normalfont File & Content\\
\midrule\endhead
\bottomrule\endlastfoot
cge\_run.gms & the driver: run \cmd{gams cge_run.gms} in the folder\\
cge\_model.gms & data, calibration, equations and the base solve\\
cge\_closure.gms & bounds, variables that do not exist, and the closure\\
cge\_solve.gms & the solve statement, for the solver and model type\\
cge\_write.gms & writes a solution\\
cge\_data.inc & the SAM, the roles of its accounts and the calibration settings\\
cge\_shock\_par.inc & parameter changes and the closure\\
cge\_shock\_var.inc & closure swaps and changes to fixed variables\\
cge\_constraints.inc & constraints, as complementarity pairs, and the scenario model \cmd{cgecon}\\
cge\_cons\_write.inc & writes the room within each constraint's bound and its marginal (from version 0.9.8)\\
cge\_bounds.inc & the attributes of \manref{cge_bound}: \cmd{.lo}, \cmd{.up}, \cmd{.scale} and \cmd{.l} (from version 0.9.8)\\
cge\_user.inc & user variables and equations (\manref{cge_variable}, \manref{cge_equation})\\
cge\_user\_write.inc & the results of the user variables\\
\end{cgetable}
The four \cmd{.gms} files of the model are copied from the package. With
\cmd{dynamic()}, the files \cmd{cge_}\argu{sim}\cmd{_par.inc} and
\cmd{cge_}\argu{sim}\cmd{_var.inc} of each period replace the two scenario
files. GAMS writes \cmd{cge_base.csv} and \cmd{cge_}\argu{name}\cmd{.csv},
with \argu{name} the value of \cmd{sim()} or \cmd{current}, or
\cmd{cge_}\argu{sim}\cmd{.csv} for each period of a dynamic run: one row per
variable and element (scenario, variable, element1, element2, value), with
the model status in the first row. The files are described in detail in
\secref{sec:gams-files}.

\cgesubsec{What is not exported}

Microsimulation and sensitivity analysis stay in Stata. The death rate is not
a GAMS parameter: it enters dynamic runs only through the population and
labour supply of the next period, which are exported; a static scenario that
changes it is exported with a note.

\cgesec{Stored results}

\cmd{cge_gams} stores the following in \cmd{r()}:
\begin{storedres}
\resgroup{Scalars}
\res{r(N_accounts)} & accounts in the exported SAM\\
\res{r(N_changes)} & parameter and fixed-variable changes, from the base or, with \cmd{dynamic()}, from the period before, in all\\
\res{r(N_constraints)} & constraints exported\\
\res{r(N_periods)} & periods exported or compared (\cmd{dynamic()})\\
\res{r(modelstat)} & GAMS model status of the scenario (\opt{}{compare}, \cmd{import()}); 1, 2, 15 and 16 mean solved\\
\res{r(base_modelstat)} & GAMS model status of the base (\opt{}{compare})\\
\res{r(maxdiff)} & largest relative difference from the simulation, or from any period (\opt{}{compare})\\
\res{r(base_maxdiff)} & largest relative difference from the Stata base (\opt{}{compare})\\
\res{r(N_compared)} & values compared, over the scenario or every period (\opt{}{compare})\\
\res{r(walras)} & value of \cmd{WALRAS} in GAMS, largest over periods with \cmd{dynamic()} (\opt{}{compare})\\
\res{r(N_imported)} & values imported (\cmd{import()})\\
\res{r(max_residual)} & largest scaled residual of the imported solution in Stata (\cmd{import()})\\
\res{r(converged)} & 1 if the imported solution solves the model in Stata (\cmd{import()})\\
\resgroup{Macros}
\res{r(name)} & name of the exported scenario or run; with \cmd{import()}, of the simulation imported\\
\res{r(closure)} & closure of the exported scenario\\
\res{r(sims)} & simulations of the periods of the run (\cmd{dynamic()})\\
\res{r(compared_with)} & the simulation or dynamic run compared (\opt{}{compare})\\
\res{r(worst_period)} & the period with the largest difference (\cmd{dynamic()} and \opt{}{compare})\\
\res{r(script)} & the script that ran GAMS (\opt{}{run})\\
\end{storedres}

\cgesec{Methods and formulas}

\textit{Scenario files.} A parameter or fixed variable whose scenario value
$\theta^1$ differs from its base value $\theta^0$ is written as
\begin{equation*}
\theta \leftarrow \frac{\theta^1}{\theta^0}\,\theta \quad (\theta^0 \ne 0),
\qquad \theta \leftarrow \theta^1 \quad (\theta^0 = 0),
\end{equation*}
so that GAMS applies the change to its own calibrated value; for a fixed
variable the right-hand side uses its level, as in
\cmd{QFS.fx('lab') = 1.02 * QFS.l('lab');}. With \cmd{dynamic()}, $\theta^0$
is the value of the period before. Numbers are written with the fewest
significant digits, 15, 16 or 17, that give back the same double.

\textit{Model types.} Under \cmd{cns} the model is the square system of
equations; under \cmd{nlp} the same equations with the constant objective
$\text{OBJ} = 0$; under \cmd{mcp} the equations, unpaired, and for each
constraint $k$ on variable $Y$ with bound $b_k$ and adjusting variable $X$,
\begin{equation*}
\text{SLK}_k \ge 0 \;\perp\; \sigma_k\,\frac{b_k - Y}{s^Y} \ge 0, \qquad
X = \bar x_k + \delta_k\,s^X\,\text{SLK}_k,
\end{equation*}
with $\sigma_k = 1$ for an upper and $-1$ for a lower bound, $\delta_k = 1$
when $X$ rises and $-1$ when it falls as the constraint binds, $s^Y$ and
$s^X$ the scales of the variables, and $\bar x_k$ the fixed level of $X$
(\secref{sec:cons-def}, \secref{sec:gams-types}). The base is solved as model
\cmd{cgebase} and a scenario with constraints as \cmd{cgecon}.

\textit{Tolerances.} The option files that GAMS writes scale the solvers'
absolute tolerances to the largest SAM cell, with
$p = \max(10^{-11},\, 10^{-13}\max_{rc} |\SAM(r,c)|)$: PATH
\cmd{convergence_tolerance} $p$; CONOPT \cmd{rtnwma} $p / \max(1, \max_{rc}
|\SAM(r,c)|)$; IPOPT \cmd{tol} $10^{-12}$ and \cmd{constr_viol_tol} $p$;
KNITRO \cmd{feastol_abs} $p$ and \cmd{feastol} $10^{-12}$.

\textit{Comparison.} For the values $g_i$ of GAMS and $x_i$ of Stata, over
the variable elements that exist in the model, except \cmd{WALRAS},
\begin{equation*}
d = \max_i \frac{|g_i - x_i|}{\max(1, |x_i|)} .
\end{equation*}

\textit{Import.} Every free variable that exists in the model, and the
adjusting variable of each constraint, takes its GAMS value; a free variable
missing from the GAMS results is an error. The imported solution is counted
as converged when the GAMS model status is 1, 2, 15 or 16 and the largest
scaled residual of the model's equations in Stata is below $10^{-8}$
(\secref{sec:gams-import}).

\cgesec{Also see}

\begin{alsosee}
\seeentry{cge_solve}{Solve the model}
\seeentry{cge_constraint}{Limits that bind in some scenarios and not in others}
\seeentry{cge_dynamic}{Solve a recursive dynamic run}
\seeentry{cge_results}{Report simulation results}
\seechapter{chap:gams}{The GAMS implementation}
\end{alsosee}
